\subsection{Gestire lo sconto su un veicolo}\label{sec:uc-sconto}
\begin{xltabular}{\textwidth}{| >{\bfseries}p{.22\textwidth} | >{\raggedright\arraybackslash}X |}
\caption{Caso d'uso: Gestire lo sconto su un veicolo}\label{tab:uc-sconto}\\
\hline
\rowcolor{gray!15} Campo & \textbf{Contenuto}\\\hline
\endfirsthead
\multicolumn{2}{c}{{\tablename\ \thetable{} -- Continua dalla pagina precedente}}\\
\hline
\rowcolor{gray!15} Campo & \textbf{Contenuto}\\\hline
\endhead
\hline
\multicolumn{2}{r}{\textit{Continua nella pagina successiva...}}\\
\endfoot
\endlastfoot
Nome & Gestire lo sconto su un veicolo\\\hline
Livello & Obiettivo Utente\\\hline
Descrizione & Il Manager applica, modifica o rimuove una promozione sul singolo veicolo. Il catalogo del Customer considera questa percentuale nel calcolo del prezzo esposto.\\\hline
Pre-condizioni & Il Manager ha effettuato l'accesso alla propria area e il veicolo è presente nel catalogo.\\\hline
Post-condizioni & La percentuale di sconto del veicolo è salvata nel database. Il successivo calcolo del preventivo considera il maggiore fra sconto del veicolo e sconto del modello, senza sommarli.\\\hline
Attori & Manager\\\hline
Test & \code{UserGoalsFunctionalTest.managerVehicleDiscount} applica il 10\%, verifica la riduzione del preventivo e controlla il ritorno al prezzo precedente dopo la rimozione. Il test attraversa servizio e persistenza H2; non esercita i comandi della GUI né il saldo dell'acquisto.\\\hline
Flusso Principale &
\begin{enumerate}[label=\arabic*.,start=1,leftmargin=*,nosep]
  \item Il Manager consulta il parco veicoli e seleziona un esemplare.
  \item Inserisce una percentuale intera compresa fra 0 e 50 e attiva il comando per applicarla.
  \item Il sistema valida la percentuale e salva lo sconto del veicolo, sostituendo quello eventualmente già presente nella stessa transazione.
  \item L'interfaccia conferma l'esito e aggiorna i dati. Alla successiva consultazione del catalogo, il preventivo viene ricalcolato considerando la promozione salvata.
\end{enumerate}\\\hline
Flussi Alternativi &
\textbf{A1 (passo 1). Nessun veicolo selezionato:} l'interfaccia richiede una selezione e non invia modifiche.\newline\newline
\textbf{A2 (passo 2). Percentuale non valida:} campo vuoto, valore non intero o fuori dall'intervallo 0--50 causano un errore; lo sconto precedente resta invariato.\newline\newline
\textbf{A3 (passo 2). Rimozione dello sconto:} il Manager sceglie il comando di rimozione; il sistema elimina la promozione individuale. L'eventuale sconto del modello resta applicabile al preventivo.\newline\newline
\textbf{A4 (passo 3). Salvataggio non riuscito:} la transazione viene annullata e l'interfaccia comunica l'errore.\\\hline
\end{xltabular}
\noindent\textbf{Limite del flusso di acquisto.} Lo sconto interviene sul preventivo di catalogo; il saldo usa invece il prezzo memorizzato nel veicolo, secondo l'anomalia descritta nella sezione~\ref{sec:pricing}.
